\documentclass{article}
\usepackage{amsmath,amssymb,amsthm}
\newtheorem{theorem}{Theorem}
\newtheorem{lemma}{Lemma}
\begin{document}
\title{A complete spectral transcript for boundary-retaining selector trees}
\author{A011, stream B03}
\maketitle
All cubes are uniform, Boolean values are signs, and entropy is in bits.
The result is a restricted-class construction, not full FEI. Independent
review and literature novelty remain unestablished.

Let $\mathcal T$ be a finite binary decision tree querying selector
coordinates $x_1,\ldots,x_a$, never querying a coordinate twice on a path.
Selector coordinates can occur at arbitrarily many nodes on different paths.
The independent data cube is $y\in\{-1,1\}^n$. Each leaf $\ell$ has a
Boolean data function $g_\ell(y)$. They can share all their variables and
Fourier supports. A leaf path specifies a coordinate subcube $C_\ell$ of
codimension $d_\ell$ and probability $p_\ell=2^{-d_\ell}$. Put
\[
 F(x,y)=g_{\ell(x)}(y),\qquad D=\sum_\ell p_\ell d_\ell.
\]
A complete leaf spectral code means a prefix code on finite witnesses
$W_\ell$, a deterministic decoder $S_\ell(W_\ell)$, and a witness probability
law whose decoded marginal is $\widehat g_\ell(S)^2$. Write $c_\ell$ for
the expected code length. Subprobability Kraft sums are allowed.

\begin{theorem}[Complete cost and physical admission]
For every tree and complete leaf spectral codes as above, there is a
complete prefix code for a joint witness $W$ of the spectrum of $F$, with
expected length
\[
 \mathbb E L(W)=2D+\sum_\ell p_\ell c_\ell.
\]
In particular $H(F)\le 2D+\sum_\ell p_\ell c_\ell$.
The represented positive event has the constructive gate cover
\[
 \{F=1\}=\bigcup_\ell \bigl(C_\ell\cap\{g_\ell=1\}\bigr).
\]
If $I_{\rm sel}(F)\ge\gamma D$ for some $\gamma>0$ and
$c_\ell\le K I(g_\ell)$ for every leaf, then the whole expected cost is
at most $\max\{2/\gamma,K\} I(F)$.
No admission claim is made for every Boolean function.
\end{theorem}

\begin{proof}
First construct the exact witness law recursively. At a node querying
$x_j$, let $G,H$ be its two actual Boolean child functions on the remaining
cube, and write $u_T=\widehat G(T)$, $v_T=\widehat H(T)$,
\[
 q_T=(u_T^2+v_T^2)/2,\qquad
 \theta_T=(u_T-v_T)^2/(4q_T)\quad(q_T>0).
\]
Choose a fair physical child $B$, sample its complete recursive witness,
and let $T$ be its decoded Fourier set. Then independently of that witness
conditional on $T$, sample a spectral bit $Z$ with probability
$\theta_T$ of being one. Decode the result as $T\cup\{j:Z=1\}$.
Zero-$q_T$ packets are never sampled. The inequality
$(u_T-v_T)^2\le2(u_T^2+v_T^2)$ makes this a valid Bernoulli law.
Its output probabilities are
\[
 q_T(1-\theta_T)=(u_T+v_T)^2/4,\qquad
 q_T\theta_T=(u_T-v_T)^2/4.
\]
These are exactly the two selector-frequency sectors of the parent.
Induction proves the decoded law is $\widehat F^2$, including all signed
cross terms. Physical path probability is $p_\ell$, since child choices
are fair regardless of the spectral choices.

Encode first the physical path to a leaf, using $d_\ell$ branch bits.
These leaf words are prefix-free. Next use the leaf's complete prefix
code. Finally append the $d_\ell$ sampled spectral bits in a prescribed
order along the known path. This is decodable and prefix-free: the physical
prefix determines the leaf dictionary and the number of final bits; the
leaf prefix determines where those final bits begin. Its Kraft sum is
\[
 \sum_\ell 2^{-d_\ell}
 \left(\sum_{w_\ell}2^{-L_\ell(w_\ell)}\right)
 \sum_{z\in\{0,1\}^{d_\ell}}2^{-d_\ell}\le1.
\]
This includes words not sampled by the true witness law. No success
conditioning or hidden normalization is required.
Since queried coordinates are distinct on a path and disjoint from the
data cube, the spectral decoder is the leaf Fourier set union the chosen
selector coordinates. Every witness has length $2d_\ell+L_\ell$.
Taking expectation gives the claimed cost. Shannon's prefix inequality
and deterministic decoding give $H(F)\le H(W)\le\mathbb E L(W)$.

The gate cover follows directly from the partition into leaf subcubes.
Conditioning on selectors gives
$I_{\rm data}(F)=\sum_\ell p_\ell I(g_\ell)$, so
$I(F)=I_{\rm sel}(F)+I_{\rm data}(F)$. Charge the $2D$ term only to
selector influence and the leaf term only to data influence. This proves
the final assertion without reusing influence.
\end{proof}

\begin{theorem}[Repeated unsigned character leaves]
Suppose $g_\ell=\chi_{A_\ell}$ with each $A_\ell$ nonempty; labels may
repeat. Define
\[
 B=\mathbb E_x\sum_{j=1}^a
       \mathbf1\{A_{\ell(x)}\ne A_{\ell(x^{\oplus j})}\}.
\]
Then $I_{\rm sel}(F)=B/2$, $I_{\rm data}(F)=\sum p_\ell|A_\ell|$,
and the full witness cost is $2D$. If $B\ge\beta D$ for a fixed
$\beta>0$, then
\[
 H(F)\le\mathbb E L(W)=2D\le(4/\beta)I_{\rm sel}(F)
 \le(4/\beta)I(F).
\]
Every gate combines a conjunction with a parity event; for singleton
$A_\ell$ every gate is an ordinary conjunction ($M=1$).
\end{theorem}
\begin{proof}
A character's spectrum is deterministic, so its leaf code is empty.
Distinct unsigned characters disagree on exactly half the data cube,
and identical ones never disagree. Average the physical selector edges.
The data identity follows as before. For any $x$, a selector coordinate
not fixed by its leaf path leaves $x$ inside the same leaf subcube when
flipped. Thus $0\le B\le D$. No spectral support disjointness is used.
\end{proof}

\begin{theorem}[An unbounded-read, vanishing-energy-per-coordinate example]
Query $x_1,\ldots,x_a$ in order, stopping at the first positive coordinate.
At stopping index $j$ output $y_1$ for odd $j$ and $y_2$ for even $j$.
The all-negative terminal leaf has label of index $a+1$ by the same rule.
For every $a\ge1$, this representation has
\[
 D=2(1-2^{-a}),\quad B\ge D/2,\quad
 I_{\rm sel}=B/2\le D/2<1,\quad I_{\rm data}=1,
 \quad \mathbb E L(W)=2D\le8I_{\rm sel}.
\]
Consequently $H(F)\le8I(F)$ uniformly, whereas $I_{\rm sel}/a\to0$.
The conjunction cover has $a+1$ gates and $x_1$ occurs in every gate.
Data labels repeat arbitrarily often.
\end{theorem}
\begin{proof}
Query $j$ is reached with probability $2^{-(j-1)}$, giving the formula
for $D$. Let $r_j$ be the chance that the continuation after query $j$
has the opposite data label from its stopping leaf. Then $r_a=1$ and
$r_j=1-r_{j+1}/2$, hence $1/2\le r_j\le1$. The label-flip probability
for coordinate $j$ is $2^{-(j-1)}r_j$: this includes both orientations
of its uniform physical edge. Therefore
$B=\sum_j2^{-(j-1)}r_j\ge D/2$. Apply the preceding theorem.
The cover is $C_\ell\cap\{y_{t(\ell)}=1\}$; every path fixes $x_1$.
\end{proof}

\paragraph{Scope and remaining line.}
The physical admission $I_{\rm sel}\ge\gamma D$, or $B\ge\beta D$,
is a hypothesis, not a proved property of arbitrary trees. The construction
does not bound the original queue reference's minimum witness cost; it
supplies a different complete reference. Expected-depth spectral coding
and decision-tree FEI bounds have prior art. This file claims a source-complete
campaign interface and its scoped applications, not literature priority.
\end{document}
